\section{Implementation Details}
\label{sec:impl-detail}

\tool is implemented with over 16,000 lines of Java code, 4,000 lines of Python, and 700 lines of shell scripts. The static analyzer uses WALA~\cite{wala}. The instrumentor and runtime agent are based on a customized version of Byteman~\cite{byteman} for dynamic instrumentation. The test runner uses a modified version of JUnit~\cite{junit}, which works with \tool's instrumentor and the runtime agent to dynamically inject faults before each test. Checkpoint/Restore In Userspace (CRIU)~\cite{criu} snapshots the JVM, speeding up the test initialization.

\para{System-Specific Error Filtering.} 
Except for filtering out 
system-agnostic functions with boolean return value such as Java collection operations 
(\Cref{sec:fia:faults-to-inject}), 
\tool filters out more system-specific errors 
according to the following criteria: 

\begin{enumerate}[wide, labelwidth=!, labelindent=\parindent, noitemsep, topsep=0pt]

\item The boolean return value of the function only involves variables declared as \texttt{final}. 
Such variables are often related to system configurations. 
Configuration errors are not within the scope of this study. 
Instead of hunting for configuration errors, 
\tool searches for fault propagation 
under valid configurations in different workloads. 

\item The return value is constant or never used in the program. 
A negation placed on those functions will not have an impact on the system.

\item The return value is calculated only from primitive-type variables. 
This happens when the function is a utility function 
used in 
implementations of classic algorithms (e.g., \texttt{isSorted()} in a sorting algorithm). 
A negation will cause incorrect calculations, which is not a fault under \tool's scope.
\end{enumerate}

Additional implementation details, 
including a dynamic call graph construction 
and a customized JUnit framework, 
can be found in \Cref{app:impl}.

